% !PRINTED none
% Abschnitt-opening / part-title page: no page number in running head. Footer furniture
% "Weber, Algebra. I." and signature "29" are ignored per standing policy.

\begin{center}
\gesperrt{Dreizehnter Abschnitt.}

\gesperrt{Die Galois'sche Theorie.}
\end{center}

\begin{center}
\rule{3cm}{0.4pt}
\end{center}

\begin{center}
\S.\,139.

\gesperrt{Der K\"orperbegriff.}
\end{center}

Ein System von Zahlen wird ein \gesperrt{Zahlk\"orper} genannt, wenn es so in sich vollendet und abgeschlossen ist, dass die vier fundamentalen Rechenoperationen (die vier Species), die \gesperrt{Addition}, die \gesperrt{Subtraction}, die \gesperrt{Multiplication} und die \gesperrt{Division}, ausgef\"uhrt mit irgend welchen Zahlen des Systems, ausgenommen die Division durch Null, immer auf Zahlen f\"uhren, die demselben System angeh\"oren. Dieser Begriff, der eine Eintheilung der Zahlenarten nach einem nat\"urlichen Gesichtspunkte giebt, ist von \gesperrt{Dedekind} eingef\"uhrt (\gesperrt{Dirichlet-Dedekind}, Vorlesungen \"uber Zahlentheorie, 2. Aufl. 1871, \S.\,159). Er ist f\"ur die Algebra von der gr\"ossten Bedeutung, und es ist nicht gleichg\"ultig, daf\"ur einen bezeichnenden und ausdrucksvollen Namen zu haben. Das Wort \gesperrt{Zahlk\"orper} ist von \gesperrt{Dedekind} nach zahlreichen Analogien gebildet, in denen das Wort K\"orper (corpus, corps) in \"ahnlicher Weise eine Vereinigung von zusammengeh\"origen Dingen, der eine gewisse Vollst\"andigkeit zukommt, bedeutet.

Der Begriff des Zahlk\"orpers kann erweitert und auf alle Gr\"ossen \"ubertragen werden, mit denen nach den Regeln der vier Species gerechnet werden kann; insbesondere also auf rationale Functionen irgend welcher Ver\"anderlichen.

Wir nennen also einen \gesperrt{Functionenk\"orper} ein System von Functionen von einer oder mehreren Ver\"anderlichen von der Beschaffenheit, dass in diesem System die Rechnungen mit den vier Species unbegrenzt ausgef\"uhrt werden k\"onnen und immer auf eine bestimmte Function desselben Systems f\"uhren (immer mit Ausnahme der Division durch Null). Die Ver\"anderlichen
%CONT>>
